% preamble.tex — pacchetti, stile e comandi condivisi del documento.
% Convenzioni per i capitoli:
%   - ogni capitolo: \chapter{Titolo}\label{cap:...}
%   - sezioni: \section{...}\label{sec:...}
%   - codice inline: \codice{...}  (monospace)
%   - percorsi file: \file{...}    (monospace)
%   - blocchi di codice: ambiente lstlisting con stile "python"
%   - figure: \includegraphics da figure/generate/ o diagrammi TikZ da
%     figure/diagrammi/diagrammi.tex (input condiviso)

\usepackage[utf8]{inputenc}
\usepackage[T1]{fontenc}
\usepackage[italian]{babel}
\usepackage[a4paper,margin=2.5cm]{geometry}
\usepackage{amsmath}
\usepackage{amssymb}
\usepackage{graphicx}
\usepackage{booktabs}
\usepackage{listings}
\usepackage[expansion=false]{microtype}
\usepackage{xcolor}
\usepackage{enumitem}
\usepackage{caption}
\usepackage{hyperref}
\usepackage{tikz}
\usetikzlibrary{positioning,arrows.meta,shapes.geometric,calc}

\hypersetup{
  colorlinks=true,
  linkcolor=blue!60!black,
  urlcolor=blue!60!black,
  citecolor=green!40!black,
  pdftitle={Intelligenza artificiale nel progetto supervised_telemedicina},
}

% ---------------------------------------------------------------------------
% Comandi custom
% ---------------------------------------------------------------------------
\newcommand{\codice}[1]{\texttt{#1}}
\newcommand{\file}[1]{\texttt{#1}}
\newcommand{\progetto}{\texttt{supervised\_telemedicina}}
\newcommand{\classe}[1]{\texttt{#1}}

% ---------------------------------------------------------------------------
% Stile listings per il codice Python
% ---------------------------------------------------------------------------
\definecolor{codebg}{RGB}{245,245,245}
\definecolor{codekw}{RGB}{0,0,139}
\definecolor{codecm}{RGB}{128,128,128}
\definecolor{codestr}{RGB}{163,21,21}
\lstdefinestyle{pythonstyle}{
  language=Python,
  backgroundcolor=\color{codebg},
  basicstyle=\ttfamily\small,
  keywordstyle=\color{codekw}\bfseries,
  commentstyle=\color{codecm}\itshape,
  stringstyle=\color{codestr},
  showstringspaces=false,
  breaklines=true,
  frame=single,
  rulecolor=\color{black!20},
  numbers=left,
  numberstyle=\tiny\color{black!50},
  numbersep=8pt,
  tabsize=4,
  captionpos=b,
  literate={à}{{\`a}}1 {è}{{\`e}}1 {é}{{\'e}}1 {ì}{{\`i}}1 {ò}{{\`o}}1 {ù}{{\`u}}1 {°}{{\textdegree}}1,
}
\lstset{style=pythonstyle}

% ---------------------------------------------------------------------------
% Ambiente per gli "esempi dal progetto" (riquadro evidenziato)
% ---------------------------------------------------------------------------
\usepackage[most]{tcolorbox}
\newenvironment{esempio}[1]{%
  \par\medskip\noindent
  \begin{tcolorbox}[colback=blue!5,colframe=blue!40!black,title={#1}]
}{%
  \end{tcolorbox}\par\medskip
}